\scanpage{31}
\mspage{\textcircled{15}}
% archive foliation: 30

\[
(25)\qquad
\diagram{\begin{tikzcd}[column sep=small,row sep=small]
&&\deleted{\searrow}&1\arrow[d]&\deleted{\swarrow}&&\\
&&&\unsure{\mathcal{G}_I^+|\mathfrak S_I^+}\times\{1,-\mathrm{id}_M\}
\arrow[dl,"\mathrm{pr}_1"']\arrow[dr,"\mathrm{pr}_2"]\arrow[dd]&&&\\
&1\arrow[r]&\unsure{\mathcal{G}_I^+|\mathfrak S_I^+}\arrow[dl,Rightarrow]\arrow[dr,"\deleted{\unsure{q}}"']&&
\{1,-\mathrm{id}_M\}\arrow[dl]\arrow[dr,Rightarrow]&1\arrow[l]&\\
1\arrow[r]&\unsure{\mathcal{G}_I^+|\mathfrak S_I^+}\arrow[dr]&&SL_\Lambda(M)\arrow[dl]\arrow[dr]\arrow[dd,"\text{cart}"']&&
\{\unsure{1},\underline{a}_Q\}\simeq\mathbb{F}_2\arrow[dl]&1\arrow[l]\\
&&\unsure{\mathcal{G}_I|\mathfrak S_I}\arrow[dl]\arrow[dr,"\mathrm{sg}"']&&
\mathfrak D_Q^+\arrow[dl]\arrow[dr]&&\\
&1&&\{\pm1\}\arrow[dl]\arrow[d]\arrow[dr]&&1&\\
&&1&1&1&&
\end{tikzcd}}
\]

\uline{Corollaire 2} L'hom. can. $GL(M)\longrightarrow\mathfrak D_Q$ induit un isomorphisme
\[
(26)\qquad
GL_\Lambda(M)_{\mathrm{ab}}
=GL_\Lambda(M)/(\unsure{\mathcal{G}_I^+|\mathfrak S_I^+}\times
\underset{\substack{\text{centre de}\\GL_\Lambda(M)}}{\underbrace{\{1,-\mathrm{id}_M\}}}
\xrightarrow{\sim}(\mathfrak D_Q)_{\mathrm{ab}}
\]
\margin{vectoriel de rang $2$ sur $\mathbb{F}_2$ et même \uline{can.} isom. à $\mathbb{F}_2\times\mathbb{F}_2$.}

\margin{\uline{NB.}
$\left\{\begin{array}{l}
\det g=\unsure{\det}(\dot g)\\
\mathrm{sg}\,g=\chi_C(\dot g)\\
\det g\cdot\mathrm{sg}\,g=\chi_D(\dot g)
\end{array}\right.$
pour $g\in GL(M)$, $\dot g$ im. can. $\in\mathfrak D_Q$.}

On peut résumer \inserted{l'essentiel de l'information de} (23) (25) dans un diagramme des sous-groupes de $GL_\Lambda(M)$, où on a posé pour abréger
\[
(27)\qquad
\unsure{Z|\mathfrak Z}=\{1,-\mathrm{id}_M\}
\left(=\mathrm{Cent}(GL_\Lambda(M))=\mathrm{Cent}(SL_\Lambda(M))\right)
\]
